% Appendix A: Detailed Proofs
% \input-able LaTeX file for Paper 5
% Conventions: a \sp b = sequential product (dot notation matching vdW)
%   C_p = Alfsen-Shultz compression (P-projection)

\section{Detailed Proofs}\label{app:proofs}

%--------------------------------------------------------------------
\subsection{Proof of \ref{ax:S4} (Orthogonality Symmetry)}
\label{app:S4-proof}

We give the full proof that the self-modeling sequential
product~\eqref{eq:corrected-product} satisfies axiom~\ref{ax:S4}: if
$\seqp{a}{b} = 0$ then $\seqp{b}{a} = 0$.  The proof depends only on the
facial geometry of spectral order unit spaces and the vanishing of the
mixing function at zero eigenvalues ($f(0, x) = 0$), making it
independent of the specific choice $f = \sqrt{\lambda_{i}\,
\lambda_{j}}$.

\begin{theorem}\label{thm:S4-full}
Let $(V, \leq, \id)$ be a finite-dimensional spectral order unit
space equipped with the sequential product~\eqref{eq:corrected-product}.
If $\seqp{a}{b} = 0$ for effects $a, b \in \eff{V}$, then
$\seqp{b}{a} = 0$.
\end{theorem}

\begin{proof}
Write the spectral decomposition $a = \sum_{i=1}^{n}
\lambda_{i}\, p_{i}$, with $\lambda_{i} \geq 0$, the $p_{i}$
mutually orthogonal projective units, and $\sum_{i} p_{i} = \id$.
The condition $\seqp{a}{b} = 0$ reads
\begin{equation}\label{eq:S4-expand}
  \sum_{i} \lambda_{i}\, \comp{p_{i}}(b)
  + \sum_{i < j} f(\lambda_{i}, \lambda_{j})\, \peirce{ij}(b)
  = 0.
\end{equation}
By the Peirce decomposition~\eqref{eq:peirce-proj}, the terms
$\comp{p_{i}}(b) \in V_{2}(p_{i})$ and
$\peirce{ij}(b) \in V_{1}(p_{i}, p_{j})$ lie in mutually
orthogonal subspaces.  In a direct sum, a sum vanishes if and only if
each summand vanishes.  Therefore~\eqref{eq:S4-expand} implies
\begin{align}
  \lambda_{i}\, \comp{p_{i}}(b) &= 0
  \quad \text{for each } i,
  \label{eq:S4-diag}\\
  f(\lambda_{i}, \lambda_{j})\, \peirce{ij}(b) &= 0
  \quad \text{for each } i < j.
  \label{eq:S4-offdiag}
\end{align}

\medskip
\noindent\textbf{Case~A: Full-rank ($\lambda_{i} > 0$ for all $i$).}
From~\eqref{eq:S4-diag}, $\comp{p_{i}}(b) = 0$ for all $i$.
From~\eqref{eq:S4-offdiag}, since $f(\lambda_{i}, \lambda_{j}) > 0$
when both arguments are positive (this holds for any mixing function
with $f > 0$ on $(0,1]^{2}$, including
$f = \sqrt{\lambda_{i}\, \lambda_{j}}$), we get
$\peirce{ij}(b) = 0$ for all $i < j$.  The completeness of the Peirce
decomposition~\eqref{eq:peirce-proj} then gives
\[
  b = \sum_{i} \comp{p_{i}}(b) + \sum_{i < j} \peirce{ij}(b) = 0.
\]
Hence $\seqp{b}{a} = \seqp{0}{a} = 0$ (by~\ref{ax:S1}, the map
$b \mapsto \seqp{b}{a}$ sends $0$ to $0$).

\medskip
\noindent\textbf{Case~B: Rank-deficient ($\lambda_{k} = 0$
for some $k$).}
Partition the indices into $I_{+} = \{i : \lambda_{i} > 0\}$ and
$I_{0} = \{i : \lambda_{i} = 0\}$, both nonempty.
From~\eqref{eq:S4-diag}: $\comp{p_{i}}(b) = 0$ for all $i \in I_{+}$.
From~\eqref{eq:S4-offdiag}: for $i, j \in I_{+}$,
$f(\lambda_{i}, \lambda_{j}) > 0$ gives
$\peirce{ij}(b) = 0$; for $i \in I_{+}$, $j \in I_{0}$,
$f(\lambda_{i}, 0) = 0$ gives no constraint on
$\peirce{ij}(b)$.

The key step is the \emph{facial absorption theorem}
(Alfsen--Shultz~\cite{AlfsenShultz2003}, Proposition~7.43):
\begin{quote}
  If $\comp{p}(b) = 0$ and $b \geq 0$, then
  $b \in \mathrm{face}(p^{\perp})$, the face generated by
  $p^{\perp} = \id - p$.
\end{quote}
Define $p_{+} = \sum_{i \in I_{+}} p_{i}$ (the support projection of
$a$).  Since $\comp{p_{i}}(b) = 0$ for each $i \in I_{+}$, and
compressions for orthogonal projective units act independently on
their respective faces, we have
$\comp{p_{+}}(b) = \sum_{i \in I_{+}} \comp{p_{i}}(b) = 0$.
By facial absorption, $b \in \mathrm{face}(p_{+}^{\perp})$, where
$p_{+}^{\perp} = \sum_{i \in I_{0}} p_{i}$.

This means $b$ is supported entirely on the zero-eigenvalue block of
$a$.  For the reverse product $\seqp{b}{a}$, write
$b = \sum_{j} \mu_{j}\, q_{j}$ (spectral decomposition of $b$).
Since $b \in \mathrm{face}(p_{+}^{\perp})$, each spectral
projector $q_{j}$ with $\mu_{j} > 0$ satisfies
$q_{j} \leq p_{+}^{\perp}$ in the face lattice.

Now consider $\seqp{b}{a}$:
\[
  \seqp{b}{a} = \sum_{j} \mu_{j}\, \comp{q_{j}}(a)
  + \sum_{j < k} f(\mu_{j}, \mu_{k})\, Q_{jk}(a),
\]
where $Q_{jk}$ denotes the Peirce projection for the $q_{j}$-basis.
For each $j$ with $\mu_{j} > 0$: since $q_{j} \leq p_{+}^{\perp}$
and $a$ is supported on $p_{+}$ (all nonzero eigenvalues correspond to
$p_{+}$), the facial orthogonality of complementary faces gives
$\comp{q_{j}}(a) = 0$.  The Peirce $1$-space terms $Q_{jk}(a)$ either
vanish by facial structure (when both $\mu_{j}, \mu_{k} > 0$,
so both $q_{j}, q_{k} \leq p_{+}^{\perp}$, and $a$ has no component
in $V_{1}(q_{j}, q_{k})$ since $a$ is supported on the complementary
face) or carry zero weight $f(\mu_{j}, 0) = 0$ (when $\mu_{k} = 0$).
Hence $\seqp{b}{a} = 0$.

\medskip
\noindent\textbf{Case $a = 0$.}  Trivial: $\seqp{0}{b} = 0$ implies
nothing, and $\seqp{b}{0} = 0$ by~\ref{ax:S1}.
\end{proof}

\begin{corollary}[$\varphi$-independence of \ref{ax:S4}]
\label{cor:S4-phi-indep}
Axiom~\ref{ax:S4} holds for the sequential
product~\eqref{eq:general-product} with \emph{any} mixing function $f$
satisfying $f(0, x) = 0$ for all $x \in [0,1]$, not only for the
faithful choice $f = \sqrt{\lambda_{i}\, \lambda_{j}}$.
\end{corollary}

\begin{proof}
The proof of \Cref{thm:S4-full} uses only: (i)~the Peirce direct sum
structure, (ii)~the vanishing $f(\lambda, 0) = 0$, and (iii)~facial
absorption.  The specific form of $f$ for positive arguments appears
only in Case~A (where $f > 0$ on $(0,1]^{2}$ suffices) and in the
$(i, j \in I_{+})$ part of Case~B (same condition).  Any mixing
function with $f(0, x) = 0$ and $f > 0$ on $(0,1]^{2}$ satisfies the
proof's requirements.
\end{proof}


%--------------------------------------------------------------------
\subsection{Proof of Local Tomography (Trace Form Non-Degeneracy)}
\label{app:lt-proof}

We give the detailed proof that the non-degeneracy of the EJA trace
form, combined with faithful self-modeling and composite minimality,
implies local tomography.

\begin{theorem}[Local tomography, detailed version]\label{thm:lt-full}
Let $V$ be a finite-dimensional simple Euclidean Jordan algebra
equipped with the self-modeling sequential
product~\eqref{eq:corrected-product}.  Let $V_{BM}$ be the minimal
composite (\Cref{ass:minimal}) of two copies $V_{B} \cong V_{M}
\cong V$ with the product-form sequential product.  Then
\[
  \dim(V_{BM}) = \dim(V_{B}) \cdot \dim(V_{M}).
\]
\end{theorem}

\begin{proof}
\textbf{Step 1: Lower bound.}
The product effects $\{a_{i} \otimes b_{j}\}$, where $\{a_{i}\}$ and
$\{b_{j}\}$ are bases for $V_{B}$ and $V_{M}$ respectively, belong to
$V_{BM}$ by construction (C1).  If these are linearly independent in
$V_{BM}$, then $\dim(V_{BM}) \geq \dim(V_{B}) \cdot \dim(V_{M})$.

\medskip
\textbf{Step 2: Non-degeneracy of the correlation form.}
Define the correlation bilinear form
\[
  B(a, b) = \tau\!\left(a \circ \varphi^{-1}(b)\right),
\]
where $\tau$ is the normalized trace functional on the EJA $V$
(Faraut--Kor\'anyi~\cite{FarautKoranyi1994}, Chapter~III), $\circ$
denotes the Jordan product, and $\varphi^{-1}: V_{M} \to V_{B}$ is
the inverse of the tracking isomorphism.

On a simple EJA, the trace form $(a, c) \mapsto \tau(a \circ c)$ is
non-degenerate (Faraut--Kor\'anyi~\cite{FarautKoranyi1994},
Proposition~III.4.2): if $\tau(a \circ c) = 0$ for all $a \in V$,
then $c = 0$.  Since $\varphi: V_{B} \to V_{M}$ is an order
isomorphism (Assumption~\ref{ass:faithful}), its inverse
$\varphi^{-1}$ is a bijection.  Therefore, if $B(a, b) = 0$ for all
$a \in V_{B}$, then $\tau(a \circ \varphi^{-1}(b)) = 0$ for all $a$,
which gives $\varphi^{-1}(b) = 0$ by non-degeneracy, hence $b = 0$.
So $B$ is non-degenerate.

\medskip
\textbf{Step 3: Linear independence of product effects.}
Suppose $\sum_{i,j} \alpha_{ij}\, a_{i} \otimes b_{j} = 0$ in
$V_{BM}$, where $\{a_{i}\}$ and $\{b_{j}\}$ are bases.  For any
state $\omega$ on $V_{BM}$, evaluating on this zero element gives
$\sum_{i,j} \alpha_{ij}\, \omega(a_{i} \otimes b_{j}) = 0$.
In particular, for product states
$\omega = \omega_{B} \otimes \omega_{M}$:
\[
  \sum_{i,j} \alpha_{ij}\, \omega_{B}(a_{i})\, \omega_{M}(b_{j}) = 0
  \quad \text{for all } \omega_{B}, \omega_{M}.
\]
The non-degeneracy of $B$ (which is defined via the Jordan product and
trace, both available on the EJA) ensures that the matrix
$(\alpha_{ij})$ must be zero: fix $j$ and vary $\omega_{B}$ to obtain
$\sum_{i} \alpha_{ij}\, \omega_{B}(a_{i}) = 0$ for all
$\omega_{B}$; since states separate points in an OUS
(Alfsen--Shultz~\cite{AlfsenShultz2003}, Theorem~1.23), this gives
$\sum_{i} \alpha_{ij}\, a_{i} = 0$ for each $j$, hence
$\alpha_{ij} = 0$ (since $\{a_{i}\}$ is a basis).

This establishes the lower bound:
$\dim(V_{BM}) \geq \dim(V_{B}) \cdot \dim(V_{M})$.

\medskip
\textbf{Step 4: Upper bound via minimality.}
The product effects $\{a_{i} \otimes b_{j}\}$ span a
$(\dim V_{B} \cdot \dim V_{M})$-dimensional subspace of $V_{BM}$.
The product-form sequential product~\eqref{eq:product-sp} preserves
this subspace: $\seqp{(a \otimes b)}{(c \otimes d)} =
(\seqp{a}{c}) \otimes (\seqp{b}{d})$, which is again a product effect.
The non-signaling constraints (C4) are satisfied by the
product-effect subspace.  By
minimality (\Cref{ass:minimal}), $V_{BM}$ contains no
structure beyond what is forced by (C1)--(C4) and the product-form SP.
Therefore $\dim(V_{BM}) \leq \dim(V_{B}) \cdot \dim(V_{M})$.

\medskip
Combining the lower and upper bounds:
$\dim(V_{BM}) = \dim(V_{B}) \cdot \dim(V_{M})$.
\end{proof}

\begin{remark}[Role of simplicity]
The simplicity assumption (\Cref{ass:simple}) enters at Step~2:
the trace form is non-degenerate on simple EJAs but not on direct
sums.  For a direct sum $V = V_{1} \oplus V_{2}$, the cross terms
$\tau(a_{1} \circ a_{2}) = 0$ for $a_{1} \in V_{1}$,
$a_{2} \in V_{2}$, so $B$ degenerates on cross-sector products.
However, the argument applies to each simple summand independently.
\end{remark}
